\chapter{简介(Introduction)}
刚体动力学常用三维向量（x,y,z）来表示。但RNEA采用六维向量进行表示，无比精妙。这样做的优点有：
\par 1.使代数式的表达更为简洁
\par 2.缩减了代数式所占据的空间（现代算法的实现工具是计算机，计算机的内存其实并不是很大。）
\par 3.简化了计算机内的运算
\section{动力学算法(Dynamics Algorithms)}
一个刚体系统的动力学一定是由一个方程，一个能描述作用在系统上的力和这些力所造成的加速度的方程，来进行表示的。
\par 那么动力学算法何为啊？动力学算法就是一个求解动力学量的数值的一个流程。
\par 我们初中时就知道$F=ma$，何为啊？所谓动力学，就是力和加速度的关系。
于是动力学便可以拆为两个部分：正动力学和逆动力学。
\par 1.正动力学(Forword Dynamics):由力求加速度
$$\ddot{\mathbf{q}}=FD(model,\mathbf{q}, \dot{\mathbf{q}},\boldsymbol{\tau})=\mathbf{H}^{-1}(\boldsymbol{\tau}-\mathbf{C})$$
\par 2.逆动力学(Inverse Dynamics):由加速度求力
$$\boldsymbol{\tau}=ID(model,\mathbf{q}, \dot{\mathbf{q}},\ddot{\mathbf{q}})=\mathbf{H}\ddot{\mathbf{q}}+\mathbf{C}$$
\par 一个刚体的运动学方程可以被表示为这样的一个普适形式：
$$\boldsymbol{\tau} = \mathbf{H}(\mathbf{q})\ddot{\mathbf{q}} + \mathbf{C}(\mathbf{q}, \dot{\mathbf{q}})$$
其中：\begin{itemize}
    \item $\boldsymbol{\tau}$是作用在系统上的外力向量
    \item $\mathbf{H}$是惯性矩阵
    \item $\mathbf{C}$是偏置力项矩阵
\end{itemize}


\section{空间向量(Spatial Vectors)}
三维空间中的刚体有六个运动自由度，但我们通常用三维矢量来表示它的动力学。因此，要表述刚体的运动方程，我们实际上必须表述两个矢量方程：
$$\mathbf{f}=m\mathbf{a}_C$$
$$\mathbf{n}_C=\mathbf{I}\dot{\mathbf{\omega}}+\mathbf{\omega} \times \mathbf{I}\mathbf{\omega}$$
第一个方程表示施加在物体上的力与其质心的线加速度之间的关系。第二个方程表示施加在物体上的力矩，即物体的质心，和它的角加速度之间的关系。（传递矩阵法为什么是四个基本方程？）
在空间矢量表示法中，我们使用6D矢量，它结合了刚体运动的线性和角度方面。这样，线加速度和角加速度组合成空间加速度矢量，力和力矩组合成空间力矢量，依此类推。使用空间矢量表示法，刚体的运动方程可以写成
$$\boldsymbol{f} = \mathbf{I}\boldsymbol{a} + \boldsymbol{v} \times^* \mathbf{I}\boldsymbol{v}$$
\begin{itemize}
    \item $\boldsymbol{f}$ (空间力/Spatial Force): 一个 6 维向量（也称为 Wrench），包含了作用在物体上的力和力矩。
    \item $\mathbf{I}$ (空间惯性矩阵/Spatial Inertia): 一个 $6 \times 6$ 的对称矩阵，包含了物体的质量、质心位置以及惯性张量。
    \item $\boldsymbol{a}$ (空间加速度/Spatial Acceleration): 描述刚体运动状态变化的 6 维向量。
    \item $\boldsymbol{a}$ (空间加速度/Spatial Acceleration): 描述刚体运动状态变化的 6 维向量。
    \item $\boldsymbol{a}$ (空间加速度/Spatial Acceleration): 描述刚体运动状态变化的 6 维向量。
    \item $\boldsymbol{v}$ (空间速度/Spatial Velocity): 包含角速度和线速度的 6 维向量（也称为 Twist）。
    \item $\times^*$ (对偶跨乘算子/Dual Cross Product): 这是空间向量代数特有的运算，用于计算力向量在旋转参考系中的变化。
\end{itemize}
这样的表示方法提供了很多优良的性质哦。比如$$\mathbf{I}_\text{new}=\mathbf{I}_1+\mathbf{I}_2$$